\section*{Chapitre \ref{ch:g9:maps-of-the-earth} --- Cartographier la Terre entière}

\begin{solution}{exo:g9:maps-of-the-earth:1}
(a) Une \omterm{def:g9:maps-of-the-earth:projection}{projection cartographique}. (b) Une \omterm{def:g9:maps-of-the-earth:contour}{courbe de niveau}. (c) Une
\omterm{def:g9:maps-of-the-earth:numerical-scale}{échelle numérique}. (d) L'\omterm{def:g9:maps-of-the-earth:contour}{équidistance}. (e) La \omterm{def:g9:maps-of-the-earth:remote-sensing}{télédétection}.
\end{solution}

\begin{solution}{exo:g9:maps-of-the-earth:2}
(a) $4 \times 25\,000 = \num{100000}$~cm $= \qty{1}{km}$. (b) $3 \times
50\,000 = \num{150000}$~cm $= \qty{1.5}{km}$. (c) $2.5 \times 1\,000\,000
= \num{2500000}$~cm $= \qty{25}{km}$. (d) $\qty{1}{km} = \num{100000}$~cm~;
$100\,000 \div 25\,000 = \qty{4}{cm}$.
\end{solution}

\begin{solution}{exo:g9:maps-of-the-earth:3}
(a) Faux~: une sphère ne peut pas être mise à plat sans déformer les
longueurs, si bien que tout \omterm{def:g3:continents-and-oceans:world-map}{planisphère} déforme quelque chose. (b) Vrai.
(c) Faux~: une \omterm{def:g2:drawing-to-map:map}{carte} équivalente conserve les surfaces, mais déforme les
formes, surtout près de ses bords. (d) Vrai.
\end{solution}

\begin{solution}{exo:g9:maps-of-the-earth:4}
(a) La pente est raide. (b) Le terrain est en pente douce. (c) Une
\omterm{def:g1:landscapes-around-us:landform}{colline}, dont le sommet se trouve à l'intérieur de la plus petite courbe.
(d) Un point du terrain n'a qu'une seule \omterm{def:g4:world-physical-map:altitude}{altitude}, il ne peut donc pas se
trouver sur deux courbes d'\omterm{def:g4:world-physical-map:altitude}{altitudes} différentes.
\end{solution}

\begin{solution}{exo:g9:maps-of-the-earth:5}
(a) À \ang{60}\,N, le \omterm{def:g1:landscapes-around-us:water}{lac} est dessiné 4 fois trop grand~: il paraît aussi
grand que $4 \times 1000 = \qty{4000}{km^2}$. À l'\omterm{def:g3:continents-and-oceans:equator}{équateur}, les surfaces ne
sont pas agrandies, donc le \omterm{def:g1:landscapes-around-us:water}{lac} dessiné avec la même taille a une surface
réelle de \qty{4000}{km^2}. (b) $4 \div 1.33 \approx 3$ fois plus grande.
(c) Les surfaces, et les distances du nord au sud, grandissent sans limite
vers les \omterm{def:g3:continents-and-oceans:poles}{pôles}~: les \omterm{def:g3:continents-and-oceans:poles}{pôles} seraient infiniment loin de l'\omterm{def:g3:continents-and-oceans:equator}{équateur}.
\end{solution}

\begin{solution}{exo:g9:maps-of-the-earth:6}
(a) \qty{50}{m}. (b) Le sommet est à \qty{544}{m}~; la courbe maîtresse
qui l'entoure est celle de \qty{500}{m}. (c) À l'ouest~: une \omterm{def:g1:landscapes-around-us:water}{rivière}
descend la pente, donc son fond est plus bas vers l'aval, à l'est.
(d) Entre les courbes d'environ 300 et \qty{400}{m}~: c'est là que la
pente est la plus raide, plus douce à la fois près du fond de la \omterm{def:g1:landscapes-around-us:landform}{vallée} et
près du sommet.
\end{solution}

\begin{solution}{exo:g9:maps-of-the-earth:7}
(a) C à environ \qty{400}{m}, D à environ \qty{530}{m}. (b) Environ
\qty{135}{m}, à environ \qty{1.3}{km} de C, au fond de la \omterm{def:g1:landscapes-around-us:landform}{vallée}.
(c) Environ 0.3 à \qty{0.4}{km}. (d) Les dénivelées sont faibles par
rapport aux distances~: sans exagération, les \omterm{def:g1:landscapes-around-us:landform}{collines} paraîtraient
presque plates.
\end{solution}

\begin{solution}{exo:g9:maps-of-the-earth:8}
(a) $2 \times 25\,000 = \num{50000}$~cm $= \qty{500}{m}$. (b) $300 - 200 =
\qty{100}{m}$. (c) $100 \div 500 = 0.2$, une pente de \qty{20}{\%}. (d) Moins raide~:
\qty{4}{cm} représentent \qty{1000}{m}, donc la pente est de $100 \div 1000 = \qty{10}{\%}$,
deux fois moins forte.
\end{solution}

\begin{solution}{exo:g9:maps-of-the-earth:9}
(a) $14 \times 18 = 252$ orbites. (b) $14 \times 103 = 1442$ minutes,
environ \qty{24}{h}~: une journée. (c) $5.5 \div 1 = 5.5$ fois. (d) Par
exemple~: un \omterm{def:g8:ice-on-earth:glacier}{glacier} qui recule, une forêt abattue, une ville qui
s'étend, une \omterm{def:g1:landscapes-around-us:water}{côte} ou une \omterm{def:g1:landscapes-around-us:water}{rivière} qui se déplace.
\end{solution}

\begin{solution}{exo:g9:maps-of-the-earth:10}
(a) $20\,200 \div 300\,000 \approx 0.067$~s $= 6.7 \times 10^{-2}$~s.
(b) $300\,000 \times 10^{-3} = \qty{300}{km}$. (c) Une infime erreur de
temps produit une énorme erreur de distance~: les horloges doivent être
extrêmement précises.
\end{solution}

\begin{solution}{exo:g9:maps-of-the-earth:11}
(a) Distance à A~: $\sqrt{3.6^2 + 4.8^2} = \sqrt{12.96 + 23.04} =
\sqrt{36} = 6$~; distance à B~: $\sqrt{6.4^2 + 4.8^2} = \sqrt{40.96 +
23.04} = \sqrt{64} = 8$. (b) $(3.6, -4.8)$, le symétrique par rapport à
la droite AB. (c) Distance de C à $(3.6, 4.8)$~: $\sqrt{3.6^2 + 5.2^2} =
\sqrt{40} \approx 6.3$~; à $(3.6, -4.8)$~: $\sqrt{3.6^2 + 14.8^2} =
\sqrt{232} \approx 15.2$. Le récepteur est en $(3.6, 4.8)$. (d) Une
distance donne un cercle (une sphère dans l'espace), deux donnent deux
points possibles~; il faut une autre distance pour choisir.
\end{solution}

\begin{solution}{exo:g9:maps-of-the-earth:12}
(a) $29 \div 2.1 \approx 14$ fois. (b) $2.1 \times 10.5 \approx
22$~millions de~\unit{km^2}. (c) Le Groenland semble couvrir environ
22~millions de~\unit{km^2}, proche des 29~millions de l'Afrique~: sur la
\omterm{def:g2:drawing-to-map:map}{carte} murale de Mercator, les deux paraissent à peu près de même taille,
alors que l'Afrique est 14 fois plus grande. (d) Une \omterm{def:g9:maps-of-the-earth:projection}{projection
équivalente}.
\end{solution}

\begin{solution}{pb:g9:maps-of-the-earth:1}
\textbf{1.} $50\,000$~cm $= \qty{500}{m}$.

\textbf{2.} $3.6 \times 500 = \qty{1800}{m} = \qty{1.8}{km}$.

\textbf{3.} B--W~: $2.5 \times 500 = \qty{1250}{m}$~; W--T~: $2.35 \times 500
\approx \qty{1180}{m}$.

\textbf{4.} Vers le nord-est.

\textbf{5.} \qty{50}{m}.

\textbf{6.} S--B descend~: $296 - 143 = \qty{153}{m}$.

\textbf{7.} B--W~: $335 - 143 = \qty{192}{m}$~; W--T~: $544 - 335 =
\qty{209}{m}$.

\textbf{8.} $209 \div 1180 \approx 0.18$, environ \qty{18}{\%}.

\textbf{9.} Une ligne droite de S à T traverserait la \omterm{def:g1:landscapes-around-us:water}{rivière} là où il n'y
a pas de pont~; le sentier descend jusqu'au pont, en B.

\textbf{10.} $1.8 + 1.25 + 1.18 \approx \qty{4.2}{km}$ à \qty{4}{km} par
heure~: environ 63 minutes~; montée $192 + 209 = \qty{401}{m}$~: environ
40 minutes de plus. Au total, environ 1~h~43~min.

\textbf{11.} Depuis T, elle descend jusqu'à B, puis monte de B jusqu'à S~: $296 -
143 = \qty{153}{m}$.

\textbf{12.} Longueur~: $1.8 + 1.25 + 1.18 \approx \textbf{4.2 km}$~;
dénivelée positive totale~: $192 + 209 = \textbf{401 m}$.
\end{solution}
